\subsection{The actual scalar native stars and the outer line correction}
\begin{statement}{P3K.62}{Proved}

Retain the actual fixed-source whole-zero \(I_1\) scalar families,
their source and Cartier line, finite-root and annular cap maps,
physical delta circle, and specified forward parameter closing.
On their protected cusp disc the full marked periods are
\[
 h:(Z,V)\mapsto(\kappa(q)qZ,\alpha(q)V),\qquad
 B:(Z,V)\mapsto(g(q)Z,q\lambda\beta(q)V),
\]
with the displayed units independent of the scalar parameter.
Use the initial physical-vertical-edge fan and its labelled native
stars \(C_1,C_2\), with the coordinates constructed below.
For an integral scalar winding \(J\), including the actual
\(\lambda=r^M\) covers and either norm sign, the single native
relative free pairs are
\[
 (m_{C_1},n_{C_1})=(J,-J),\qquad
 (m_{C_2},n_{C_2})=(0,0).
 \tag{P3K.62.0}
\]
The actual first \(C_1\) flop changes the oriented central return
by \(-J\) in \(\Theta_7\cong\mathbb Z/28\). The actual first
\(C_2\) flop preserves that return and its integer characteristic
index. In particular the initial and reflected fan scalar
indices used in P3K.58 are equal by this explicit \(C_2\) map.

For \(C_1\), the two outer positive-Hopf normal coordinate maps
have determinant windings \(0\) and \(-2J\). This is the missing
outer line marking in P3K.25's earlier zero-correction argument.
For \(J\ne0\), the actual labelled delta-equivariant simultaneous
three-curve plumbing product invoked in P3K.29 is obstructed by
its full normal determinant square. The local topology and
single-star frame maps remain the ones used below. This statement
does not assign a universal all-fan formula.
\end{statement}
\begin{proof}
\par

\medskip\noindent\textbf{SCN.1. The actual primitive scalar seam}

\par\nobreak

Keep the physical line character \(V\) and the original positive
\(\lambda\)-loop of K22/K25. The full periods are
\[
 h:(Z,V,q)\mapsto(\kappa(q)qZ,\alpha(q)V,q),\qquad
 B:(Z,V,q)\mapsto(g(q)Z,q\lambda\beta(q)V,q),
\]
where all displayed units are independent of the scalar parameter.
Put \(\eta=\lambda\beta\). The primitive cap-to-ordinary gluing of
K10/K22 has the literal forward parameter arrow
\[
                  R_{\rm ord}:(Z,V,q)\mapsto(Z,q^{-1}V,q).
                                                        \tag{SCN.1}
\]
Its cap arrow is the holomorphic unit recorded in K10; the exact
overlap square there is retained unchanged. Reindexing the shifted
cone by the actual period \(B\) gives
\[
                    E_\lambda:(Z,V)\mapsto(gZ,\eta V).
                                                        \tag{SCN.2}
\]
This is the K25.6 endpoint in global torus coordinates; in its
local square coordinates \((Z_{\rm loc},V_{\rm loc})=(Z,ZV)\)
it is \((gZ_{\rm loc},g\eta V_{\rm loc})\).

On a cut scalar interval \(\tau\in[0,1]\), let
\(\lambda=\lambda_*e^{2\pi i\tau}\). Retain the actual logarithms
of \(g,\eta\) on the \(q\)-disc times this interval, and take a
cutoff \(\chi=0\) near the first end, \(\chi=1\) near the last.
Use
\[
 C^{(1)}_\tau:(Z,V)\mapsto(AZ,DV),\qquad
 A=e^{-\chi\log g},\quad D=e^{-\chi\log\eta}.         \tag{SCN.3}
\]
On the terminal collar \(C^{(1)}=E_\lambda^{-1}\), so the
entire seam identity \(C_{\rm end}^{(1)}E_\lambda=1\) holds
with every derivative. These maps preserve \(q\), commute with
both periods and the physical circle, and have the two actual
chart expressions
\[
 (a,b,c)\mapsto((AD)^{-1}a,Ab,Dc),\qquad
 (u,v,w)\mapsto(ADu,D^{-1}v,A^{-1}w).
                                                        \tag{SCN.4}
\]
Substitution in each of DZ.8's three transition rows preserves
every coefficient. Thus these are actual closed vertex charts,
with the original units and forward \(q^{-1}\) arrow.

The doubled seam is specified as the literal pullback
\[
 C^{(2)}_{j+\tau}=C^{(1)}_\tau E_{\lambda(\tau)}^{-j},
                       \qquad j=0,1.                \tag{SCN.5}
\]
Its middle and terminal collar identities are respectively
\(E^{-1}\) and \(E^{-2}\). It has not been independently chosen
from the doubled endpoint.

\par

\medskip\noindent\textbf{SCN.2. The complete native C1 product and Cox lift}

\par\nobreak

In these closed charts the horizontal native star is given by
\[
 (z,n_0,n_1)=(c,a,b),\qquad
 (z',n'_0,n'_1)=(\eta v,g w,(g\eta)^{-1}u).
                                                        \tag{SCN.6}
\]
The full \(C_1\) transition row is
\[
              (u,v,w)=(g\eta bc,\eta^{-1}c^{-1},g^{-1}ac).
\]
Hence (SCN.6) satisfies exactly
\[
                         z'=z^{-1},\qquad n'_i=zn_i.
                                                        \tag{SCN.7}
\]
This is a single native star product with all parameter units
included. On its first chart (SCN.4) acts by
\[
                 (z,n_0,n_1)\mapsto
                    (Dz,(AD)^{-1}n_0,An_1).
                                                        \tag{SCN.8}
\]
On the second it is
\((z',n'_0,n'_1)\mapsto(D^{-1}z',A^{-1}n'_0,ADn'_1)\);
the equality (SCN.7) is retained literally.

For Cox coordinates with auxiliary weights \((1,1,-1,-1)\),
put \(z=a_1/a_0,n_i=a_0c_i\). The lift of (SCN.8) is
\[
 (a_0,a_1,c_0,c_1)\mapsto
       (a_0,D a_1,(AD)^{-1}c_0,Ac_1).                \tag{SCN.9}
\]
It preserves \(q=a_0a_1c_0c_1\), so its inverse has the same
\(q\) and inverse multipliers even for the full \(q\)-dependent
units. It extends through both small resolutions. In this order
the physical circle weights are \((0,1,-1,0)\). The angular
factor of (SCN.9) is the corresponding physical-circle factor;
the inverse direction from K25.7 is retained by (SCN.3).
The native product pulled back under the double parameter cover
is exactly the one using (SCN.5), including (SCN.6).

\par

\medskip\noindent\textbf{SCN.3. The doubled ordinary normal maps and trace degrees}

\par\nobreak

The raw period-unit winding is
\(\left(\begin{smallmatrix}0&0\\0&1\end{smallmatrix}\right)\)
for one positive primitive loop. For two turns, with
\(\lambda=\lambda_*e^{2i\theta}\), the actual based normal
ratios in the second-vertex product
\((u,\eta v,\kappa w)\) are exactly
\[
 K_0=K_2=\operatorname{diag}(e^{2i\theta},1),
 \qquad K_1=\operatorname{diag}(1,e^{2i\theta}).       \tag{SCN.10}
\]
All other units cancel against their identical reference values.
The original higher normal terms are carried by the same
normal-rescaling and cutoff-flow construction as DZ.2. Thus the
closed ordinary product uses
\[
 t=1,\qquad l=(2,0,2),\qquad
 T_i(Q)=H_1 QH_{l_i-1}.
\]
The physical weights remain \(w=(0,1,-1)\),
\(d=(-2,-1,-1)\). Substitution in the literal orbit map DZ.16
and its closed stable derivative gives
\[
                  I=(-2,2,-2),\qquad W=(2,0,0).
                                                        \tag{SCN.11}
\]
The actual capped graph-loop prism and the ambient-frame
comparison then give on \(C_1\)
\[
 V_{11}=2,\qquad\delta_1=-2,\qquad
       2n_1=\delta_1-V_{11}=-4,\qquad m_1=0.
                                                        \tag{SCN.12}
\]
The native-to-ordinary boundary arrow in this star is
\[
                    q_H\mapsto
                   (U H_1U^{-1})q_H H_{-1}.
                                                        \tag{SCN.13}
\]
Its Hopf-oriented left degree is \(+2\), and its right degree
is zero. Extending its inverse by OR.2's actual cylinder map
gives the pure-left endpoint \(L^{-2}\). The doubled native
pair obtained from (SCN.12) is therefore \((2,-2)\).

The product (SCN.5)--(SCN.7) is the literal pullback of the
single native product, so the same fixed-boundary connection
argument of DZ.4 identifies its relative return with \(f^2\).
Its two free coordinates consequently give the single native
pair
\[
                               (m(f),n(f))=(1,-1).
                                                        \tag{SCN.14}
\]
This uses the free coordinates of the actual relative square;
it does not divide a seven-sphere class.

\par

\medskip\noindent\textbf{SCN.5. The actual outer normal marking omitted by K25}

\par\nobreak

The following calculation identifies the outer normal-line mismatch.

Keep every unit of SCN.1. The native fixed-curve \(C_0\) charts
have first base \(b\), normal pair \((a,c)\), and second base
\(\kappa w\), normal pair
\[
             \bigl(\alpha v,(\kappa\alpha)^{-1}u\bigr).
                                                        \tag{SCN.15}
\]
The negative physical weight is the first normal coordinate
in both ordered pairs. In the actual \(C_1\) native chart
\((z,n_0,n_1)=(c,a,b)\), the fixed branch is
\(z=n_0=0\), with base \(n_1=b\) and normal pair \((a,c)\).
Thus its first endpoint agrees literally with (SCN.15).

In its second chart
\[
              (z',n'_0,n'_1)
                   =(\eta v,g w,(g\eta)^{-1}u),
\]
the fixed branch is \(z'=n'_1=0\), with base \(n'_0=g w\)
and ordered negative/positive normal pair
\[
                    \bigl(\eta v,(g\eta)^{-1}u\bigr).
                                                        \tag{SCN.16}
\]
Consequently the actual coordinate map from the native
\(C_0\) normal pair to this \(C_1\) pair at the second branch
has base factor \(g/\kappa\) and normal matrix
\[
       R=\operatorname{diag}
          \left(\frac{\eta}{\alpha},
                       \frac{\kappa\alpha}{g\eta}\right).
                                                        \tag{SCN.17}
\]
This is an equality of the displayed coordinates, with no
change of period or omission of a bounded multiplier.
The first branch has the identity normal matrix.

Divide (SCN.17) by its reference value at \(\lambda_*\).
The units \(\kappa,\alpha,g,\beta\) are independent of the
scalar parameter. Hence on the positive fixed-radius loop,
\[
        R R_*^{-1}
           =\operatorname{diag}(e^{i\theta},e^{-i\theta}).
                                                        \tag{SCN.18}
\]
The physical positive-weight complex normal structure
conjugates the first coordinate. Thus the actual Hopf
spinor coordinate map in (SCN.18) is
\[
                         e^{-i\theta}I_2.             \tag{SCN.19}
\]
At the other endpoint it is \(I_2\). Its determinant winds
\(-2\) at the second endpoint and \(0\) at the first.
For \(\lambda=r^m\) these windings are \(-2m\) and \(0\).
They have been calculated from the actual arrows, not
selected from SCN.14 or from a desired flop sign.

K25.7 and K25.8 compare the two \emph{inner} charged-cap maps.
Their identical physical-circle gauge can be removed
there by one cap frame. Equations (SCN.15)--(SCN.19)
compare those native master charts to the two \emph{outer}
fixed-curve caps. They give different endpoint gauges.
Multiplying both by a common parameter scalar leaves
their ratio \(e^{-i\theta}\) unchanged. Therefore the
inner common-gauge calculation alone does not establish
the claimed pullback of the outer normal-sphere line.
The precise unsupported sentence is K25's assertion,
after identifying the four fixed-boundary charges,
that its parameter circle is a separate factor and
consequently \(b^2=0\).

\par

\medskip\noindent\textbf{SCN.6. The integral Cech class with an explicit orientation}

\par\nobreak

Here is the line class of the actual mismatch, separately
from the orientation of the full characteristic bordism.
Let a spatial circle be the union of a master interval
from endpoint \(0\) to endpoint \(1\) and a cap interval
traversed from endpoint \(1\) back to endpoint \(0\).
Orient its parameter circle by increasing \(\theta\).
Use the coordinate arrows from the master frame to the
cap frame
\[
        y_{\rm cap}=g_j(\theta)y_{\rm master},\qquad
                 g_0=1,\quad g_1=e^{-im\theta}.
                                                        \tag{SCN.20}
\]
These are (SCN.19), with the indicated arrow and ordering.
They determine a line \(L_m\) on the resulting two-torus.
In this explicitly prescribed spatial orientation,
\[
       \langle c_1(L_m),[S^1_{\rm spatial}
                              \times S^1_\theta]\rangle=-m.
                                                        \tag{SCN.21}
\]
One direct verification keeps the connection arrow.
Set \(A_{\rm master}=0\). For the convention
\(\nabla=d+A\), equation (SCN.20) forces
\(A_{\rm cap}=-dg_j\,g_j^{-1}\), hence \(0\) at endpoint
\(0\) and \(im\,d\theta\) at endpoint \(1\).
Parametrize the cap in the temporary direction \(0\to1\)
and interpolate with \(A_{\rm cap}=im\rho(s)d\theta\),
where \(\rho\) is constant \(0\) and \(1\) on the
respective collars. Its curvature integral in that
temporary direction is \(m\). The oriented spatial
circle traverses the cap in the opposite direction,
giving (SCN.21). This also gives the integral Cech
calculation without differentiating a multivalued log.
The positive Hopf normal two-plane has the scalar
transition (SCN.20) on both summands; its determinant
crossing class is therefore \(2c_1(L_m)\), with
evaluation \(-2m\).

This calculation identifies a nonzero relative line
class and fixes its sign in the displayed two-torus
orientation. We have not yet identified that oriented
two-torus with the oriented four-dimensional fixed
base in K25's full bordism, nor computed its sphere
class or the resulting value of \(b^2\). No numerical
characteristic correction is assigned by SCN.21.

\par

\medskip\noindent\textbf{SCN.7. The simultaneous-product premise also fails directly}

\par\nobreak

Apply the complete determinant square PO.5--PO.11
of K47 to the actual scalar charts (SCN.4), retaining
their labelled vertices and all three labelled curves.
The pointed \(C_2\) and \(C_1\) products force
\[
                 a_*+u_*=0,\qquad c_*+v_*=m.
\]
The arbitrary invertible positive-weight complex
normal derivative along \(C_0\) forces
\[
      (-v_*+u_*)+\operatorname{wind}\kappa
            +2\operatorname{wind}\alpha-(-a_*+c_*)=0.
\]
Substitution yields
\(\operatorname{wind}\kappa+2\operatorname{wind}\alpha
 -\operatorname{wind}\eta=-m=0\), which is impossible
for \(m\ne0\). The argument allows full off-diagonal
\(GL_2(\mathbb C)\) normal matrices; it does not assume
diagonality away from the vertices. Thus the actual
marked delta-equivariant simultaneous plumbing product
asserted in K29 Section 3 does not exist for this
positive scalar loop. The separate native \(C_1\)
product of SCN.6--SCN.9 and the ordinary quaternion
product are not excluded by this determinant square.

The exact local maps and integral mismatch are recorded
by \path{code/part3/scalar_c1_outer_marking.py} and its output.
The dependency notice is kept separately in
\path{notes/part3_scalar_consistency_dependencies.md}.

\par

\medskip\noindent\textbf{SCN.8. The complete native C2 comparison}

\par\nobreak

The diagonal \(C_2\) star in these same seam-closed
charts has the actual coordinates
\[
        (z,n_0,n_1)=(a,b,c),\qquad
        (z',n'_0,n'_1)=(u,v,w).
\]
Its full transition is
\[
       (u,v,w)=(a^{-1},ab,ac),\qquad
                 z'=z^{-1},\quad n'_i=zn_i.           \tag{SCN.22}
\]
There is no omitted unit or parameter factor in this
transition. The seam (SCN.4) acts on its first chart by
\[
       (z,n_0,n_1)\longmapsto
                    ((AD)^{-1}z,An_0,Dn_1)
\]
and on the second by
\((z',n'_0,n'_1)\mapsto(ADz',D^{-1}n'_0,A^{-1}n'_1)\).
With \(z=a_1/a_0,\ n_i=a_0c_i\), its full Cox lift is
\[
   (a_0,a_1,c_0,c_1)\longmapsto
                (a_0,(AD)^{-1}a_1,Ac_0,Dc_1).        \tag{SCN.23}
\]
It preserves \(q=a_0a_1c_0c_1\), the auxiliary action,
the physical circle and every actual period translate.
It extends over both small resolutions and has the
displayed inverse at the same \(q\).

For two primitive scalar turns, SCN.11 and the actual
prisms DZ.19, with \(w_2=-1,w_0=0\), give
\[
                    V_{22}=2,\qquad \delta_2=2.
\]
The local ambient-to-product frame comparison therefore
gives the ordinary relative pair \((-2,0)\).
The native-to-ordinary boundary map is literally
\[
                q_H\longmapsto
                 (UH_1U^{-1})q_HH_1.                 \tag{SCN.24}
\]
Its left degree is \(+2\) and its right degree is zero.
The inverse extends by the explicit cylinder retraction
of OR.2, with endpoint \(L^{-2}\); its coefficient maps
extend over \(D^3\times D^4\). Thus the native doubled
pair is \((0,0)\).

Equations (SCN.5), (SCN.22) and (SCN.23) are the literal
pullback of one native \(C_2\) product. Consequently
the doubled relative return is its actual square.
The free coordinates of its single relative return
are \((0,0)\), since multiplication by two is injective
on these two integral coordinates. The actual
single-star Cox conjugation changes the central
seven-sphere class by its right integer, which is
zero. This establishes equality of the initial and
reflected scalar return classes without division in
\(\Theta_7\).

The same maps retain every integral scalar winding
\(J\). On the double pullback use
\[
       t=J,\quad l=(2J,0,2J),\quad
       I=(-2J,2J,-2J),\quad W=(2J,0,0).
                                                        \tag{SCN.25}
\]
The ordinary \(C_2\) pair is \((-2J,0)\), its native
pair is \((0,0)\), and the single native pair is
again \((0,0)\). This includes negative winding and
zero winding, with the same physical circle.

For the canonical cap families the ordinary closing
may be \(q^e\) with either sign of \(e\), rather than
the primitive cap-to-ordinary \(q^{-1}\) of SCN.1.
The exact reindexing is then \(B^{-e}\), and its
bounded map is
\((Z,V)\mapsto(g^{-e}Z,\eta^{-e}V)\).
Use the cut seam
\[
       (Z,V)\mapsto
         \bigl(e^{e\chi\log g}Z,e^{e\chi\log\eta}V\bigr).
                                                        \tag{SCN.26}
\]
It obeys \(C_{\rm end}E=1\) with all derivatives.
Its double is still the literal pullback (SCN.5).
Thus the full source and cap closing squares already
proved in K26/K28/K32 are carried unchanged; no norm
sign is reversed. Parameter-independent bounded
closing units are carried through their fixed
logarithms. The unit winding in the actual period
matrix is still \((0,0;0,J)\), so (SCN.25) applies.
The scalar equality needed by K58 for its initial
and reflected fans follows from these actual maps.

\par

\medskip\noindent\textbf{SCN.9. The C2 characteristic bordism has no outer mismatch}

\par\nobreak

There is also a direct narrow repair of the integer
characteristic comparison. At the first fixed branch
of the \(C_2\) native star, its fixed base and normal
pair are exactly \(b\) and \((a,c)\). At the second,
they are \(w\) and \((v,u)\). Relative to the native
\(C_0\) coordinates (SCN.15), the second arrow is
\[
 \text{base factor }\kappa^{-1},\qquad
 \text{normal matrix }\operatorname{diag}
                           (\alpha^{-1},\kappa\alpha).
                                                        \tag{SCN.27}
\]
Both this map and its first-branch identity are
independent of the scalar parameter. After carrying
their fixed logarithms, the outer normal maps are
constant in that parameter.

In the \(C_2\) Cox coordinates of (SCN.23), the
parameter factor has physical weights
\((0,-1,0,+1)\). The two inner charged boundary maps
therefore have one identical scalar gauge, as in
the actual master and joint-cap maps of K17.
Use that same scalar on every Cox chart and on the
joint cap. It commutes with their actual transition
functions. The outer identifications (SCN.27) are
constant in the parameter, so this gauge matches
the two outer conifold products as well. Unlike
SCN.17, there is no endpoint-dependent parameter
transition left on the fixed-cobordism closure.

The line on the closed five-dimensional orbit
correction is consequently the pullback of its
parameter-zero line. The same is true of its
tangent bundle under the literal moment-chart
maps and constant bounded-unit isotopies. Its
six-dimensional numbers \(x^3,p_1x\) vanish.
On the additional normal-sphere boundary the
closed four-dimensional fixed base is the
oriented \(S^2\) mapping torus crossed with the
parameter circle, with all its normal and line
transitions pulled back from the three-dimensional
factor. Thus its actual class \(b\) has \(b^2=0\),
and the base tangent Pontryagin number is zero.
The sphere-bundle identity
\[
       \frac{4\int x^3-\int p_1x}{24}
                 =\frac{3\int b^2-\int p_1(T\bar X)}{24}
\]
gives zero for this retained boundary. This proves
equality of the two integer characteristic indices
for this particular \(C_2\) comparison. It proves
no equality for the \(C_1\) comparison, whose
actual outer map is (SCN.17).

\par

\medskip\noindent\textbf{SCN.10. The full actual single-star flop and its central coordinate}

\par\nobreak

For each of SCN.9 and SCN.23, transport the Hermitian
metric by the displayed invertible Cox map and then
apply its positive radial normalization on the common
link. Its inverse is literal at the unchanged height
\(q\), so this construction retains all bounded
moduli and all parameter collars. The unrestricted
four-weight identity FF.1/SF.8 intertwines the two
resolution links with
\[
  \tau(h,q_H)=(q_H^{-1}h q_H,q_H),\qquad
  T(y,q_H)=(q_H^{-1}yq_H,q_H).                         \tag{SCN.28}
\]
Both are independent of the parameter in the native
products. The normal-rescaling homotopy carries the
full \(q\)-dependent unit germs, rather than only
their central linear terms. The coordinate maps
commute with every actual period translate, so one
free representative extends to the full periodic
orbit of disjoint flop curves. The source and cap
maps off these matched charts stay the original ones.

The first and second small resolutions in this arrow
use the same \(q_H\). The second complex-normal
convention is restored by the fixed anti-quaternion
map of K23/K31; its oriented central coordinate is
unchanged. In these actual conventions the relative
map is the literal conjugation
\[
 T(L^\mu R^\nu\mathcal C^s)T^{-1}
          =L^{-\nu}R^{\mu+2\nu}\mathcal C^{s+\nu}.
                                                        \tag{SCN.29}
\]
The unrestricted four-weight calculation is needed
here because each seam acts on all four Cox
coordinates, not only one pair of normal weights.

For scalar winding \(J\), SCN.25 gives ordinary
\(C_1\) pair \((0,-2J)\) on the double cover. Its
native boundary change has left degree \(2J\) and
right degree zero, yielding native double pair
\((2J,-2J)\). The connection pulled back from the
single native product gives the actual relative
square, exactly as in SCN.8. Its single free pair
is therefore \((J,-J)\). Equation (SCN.29) changes
the actual central return by \(-J\) for this first
\(C_1\) flop. For the first \(C_2\) flop the single
free pair is \((0,0)\), and the central change is
zero. No central class is divided by two.

Root has passed the complete SCN.1--9 audit, and
fourfold has independently passed SCN.8--9 including
the actual outer/inner line comparison. The chart
and integer identities are executed in the retained
39-check certificate. This result repairs only
these two actual first-flop comparisons. It supplies
no simultaneous delta-equivariant plumbing product
or universal all-fan formula.
\end{proof}
